\begin{abstract}
We study whether neural thicket ensembles built by perturbing a single pretrained checkpoint improve robustness to adversarial attacks and out-of-distribution (OOD) inputs. Our main question is whether explicit perturbation-direction design helps beyond simple random perturbations. We train a CIFAR-10 ResNet-18 baseline, generate $M=6$ experts with matched perturbation scale ($\sigma=0.015$), recalibrate batch normalization statistics, and compare three ensemble variants: random Gaussian directions, orthogonalized directions, and adversarially selected directions. We evaluate clean accuracy, expected calibration error, negative log-likelihood, projected-gradient-descent (PGD) robust accuracy at $\epsilon\in\{2/255,4/255,8/255\}$, OOD detection on SVHN, and diversity diagnostics.

Random perturbation ensembling improves strong-attack robustness over a single model at $\epsilon=8/255$ by $+0.45$ percentage points (bootstrap 95\% CI $[+0.15,+0.80]$, $p=0.004$), but reduces clean accuracy by $-1.29$ points (95\% CI $[-1.72,-0.86]$, $p<0.001$). Orthogonal perturbations are statistically indistinguishable from random for PGD@8/255 ($+0.10$ points, 95\% CI $[-0.20,+0.40]$, $p=0.599$). Adversarial-direction perturbations produce much higher disagreement but worse calibration and weaker robustness than random. These results show a measurable but modest robustness gain from perturbation ensembling, while indicating that unconstrained diversity is not sufficient for robust generalization.
\end{abstract}
